% 本文件由 LaTeX 排版 Agent 根据有效 translation.json 自动生成。
% author=Ambrose; work_id=3408; title=论寡妇

\chapter{论寡妇}

% structure: p0001 source=h1 target=omitted reason=page-title-already-rendered-by-parent

% structure: p0002 source=h2 target=section reason=relative-heading-rank; base=h2

\section{导言}

\Person{圣安博罗削}在本文开头告知我们，他在讲了有关童贞女的事之后，受\Person{圣保禄}的榜样所感，继续论及寡妇这一主题。然而，在他自己的教区内还有另一件事，对他个人触动颇深，促使他立即着手处理这一主题。有一位寡妇，膝下有几个女儿，其中一些已经出嫁，另一些也已到了适婚年龄，她开始考虑自己再婚。\Person{圣安博罗削}部分是为她着想，部分是为了不让人以为他以任何方式建议了这一步，因此发表了以下论述。\par

他首先确认，守寡的圣召与童贞的圣召非常接近，并远远高于已婚状态。他通过\Person{圣保禄}的见证及其对真寡妇的描述予以证明；同时也从旧约和新约中引用许多例子加以证实。在提及\Person{圣伯多禄}的岳母后，他更为直接地转向那位促使他写作的寡妇，尽管他避免提及她的名字，指出她为自己再婚所提出的种种理由是多么空虚和不足。他表示，婚姻的纽带确实是神圣而美好的，已婚者与独身者犹如教会田野中各式各样的花朵。然而，所产的麦粒多于百合，已婚者多于守贞者。他指出，守寡受到拜偶像者的鄙视，正因如此，它理当受到基督徒的尊敬。\Person{圣安博罗削}并不谴责再婚，但他将守寡置于再婚之上，因为他的目标是引导那些托付给他慈爱照顾的人，迈向尽可能高的成全程度。\par

这部论著是在\WorkTitle{论童贞女}之后不久写成的，即在公元377年后不久。\par

% structure: p0006 source=h2 target=section reason=relative-heading-rank; base=h2

\section{第一章}

在写过有关童贞女的著作之后，似乎也有必要谈谈寡妇，因为\Person{圣保禄}将这两类人连在一起，而后者犹如前者的教师，并且远胜于那些已婚者。\Person{厄里亚}曾被派遣到一位寡妇那里，这是一项极大的尊荣；然而，寡妇若不效法\Place{匝尔法特}寡妇的美德，特别是她的好客精神，便不会像她那样受人尊敬。在这里，他谴责了人们的贪婪，他们因贪婪而丧失了天主的恩许。\par

1. 既然我已用了三卷书论述童贞女的尊荣，弟兄们，如今按顺序，应当有一部论寡妇的著作紧随其后；因为我不可不加尊荣地略过她们，也不可将她们与应归于童贞女的称许分开，因为\Person{圣保禄}的声音已将她们与童贞女连在一起，如经上所记：\BibleQuote{没有丈夫的妇女和童女，所挂虑的是主的事，一心使身体圣洁；} \BibleRef{格前7:34} 因为在某种意义上，守寡的榜样加强了童贞的教导。那些保持婚床未受玷污的人为童贞女作证，表明贞洁应为天主而持守。戒绝一度曾是乐事的婚姻，比起对婚姻的欢愉始终一无所知，其德行几乎并不逊色。她们在两方面都是坚强的：她们对婚姻不感遗憾，因为她们持守了婚姻的信实，而且不使自己卷入婚姻的快乐之中，免得显得软弱而不能自理。\par

2. 然而，在这特殊的美德中也蕴涵着自由的奖赏。因为：\BibleQuote{妻子当丈夫活着的时候，是受束缚的；丈夫若死了，她便可以自由，随意嫁给谁，但必须是嫁于主内的人。然而，按我的意见，如果她安于现状，更是有福；我想我也有天主的圣神。} \BibleRef{格前7:39} 显然，\Person{圣保禄}表达了其中的差别，他说一个是受束缚的，而另一个是更有福的，并且他声称，这并非仅出于自己的判断，而是来自天主圣神的灌注；因此这一决定应被视为属天的，而非属人的。\par

3. 在整个人类陷入饥荒、\Person{厄里亚}被派往寡妇那里的时候，这事又有何教训？且看，如何为每一位预定了她独有的恩宠。有一位天使被派遣到\Person{童贞玛利亚}那里\BibleRef{路1:26}，一位先知被派遣到寡妇那里。还要注意，在一处是\Person{加俾额尔}，另一处是\Person{厄里叟}。可见，天使和先知中最为杰出的领袖被拣选出来。然而，守寡本身并非值得称许，除非加上守寡的美德。事实上，曾有众多寡妇，但只有一位被置于所有人之上，在这件事上，与其说是阻止他人追求，不如说是以美德的榜样激励她们。\par

4. 开头所说的话使人侧耳倾听，尽管理解之纯朴本身足以吸引寡妇追求美德的榜样；因为每人之所以出众，似乎不是由于她的身份，而是由于她的功绩，并且好客的恩宠并未被天主所忽略，正如祂在福音中亲自讲述的，以永生的极丰厚的赏报回报一杯凉水，并以源源不断、取之不尽的富裕补偿那少量的面和油。因为如果一个外教人尚且说，朋友的一切财物都应当公用，那么亲属的财物更应当何等公用！因为我们彼此是亲属，被连结成同一个身体。\par

5. 但我们并不受任何规定的款待界限所约束。因为你们为何认为这世间之物是私有财产，而这世界本是共有的呢？或者你们为何认为地上的果实是私有的，而大地本身却是共有财产呢？祂说：\BibleQuote{看，天空的飞鸟，它们不播种，也不收割。}\BibleRef{玛 6:26} 因为对于没有任何私有财产的人而言，什么也不缺乏，而天主，祂自己言语的主宰，知道如何持守祂的恩许。再说，飞鸟也不积蓄，却得饱食，因为我们的天父喂养它们。但我们却将普遍教训的警告转而谋求私利，天主说：\BibleQuote{地上所有结种子的果树，和一切结种子的菜蔬，都给你们作食物；至于地上的各种野兽，天空中的各种飞鸟，以及在陆地上爬行有生命的各种动物，我把一切青草给它们作食物。}\BibleRef{创 1:29} 通过积蓄我们反而陷入匮乏，通过积蓄我们变得空虚。因为我们若不遵守这话语，便不能企望那恩许。关注款待的诫命，乐于施予异乡人，这对我们也是有益的，因为我们也是这世上的异乡人。\par

6. 然而那位寡妇是何等圣善！她虽被极度的饥饿所困，仍持守对天主的应有敬畏，不单为自己取用食物，而是与她的儿子分享，好使自己不致比爱子活得更久。亲情固然责任重大，但宗教虔诚所带来的回报更为丰厚。因为，正如无人应被置于她的儿子之前，天主的先知也当被置于她的儿子和她的性命之前。我们应当相信，她给予他的不是少许食粮，而是她生命的全部依靠，她自己则一无所留。她是如此好客，竟将一切付出；她是如此充满信德，竟即刻相信了。\par

% structure: p0014 source=h2 target=section reason=relative-heading-rank; base=h2

\section{第二章}

宗徒关于\OriginalTerm{真正的寡妇}{widow indeed}的诫命已被列出，例如：她应教养子女，孝敬父母，渴望悦乐天主，表现得无可指摘，展示成熟的功绩，只作过一个丈夫的妻子。然而，圣盎博罗削指出，圣保禄并未谴责再婚，并补充说，寡妇必须在众人前享有良善的德性名声。还应说明为何应避开年轻寡妇，以及所谓\Quote{与其欲火攻心，倒不如嫁娶为妙}的含义。接着，圣盎博罗削又论述了寡妇的尊严，这尊严可由如下事实显明：任何对她们的伤害，都会招致天主的义怒。\par

7. 所以，寡妇并非仅由身体的节制而被分别出来，而是以德性而与众不同。对此，我不发命令，而是宗徒。我并非唯一尊敬寡妇的人，\OriginalTerm{外邦人的导师}{Doctor of the Gentiles}首先这样做了，他说：\BibleQuote{要敬重那些真正守寡的寡妇。若寡妇有儿女或孙儿孙女，她应先学习管理家务，报恩于父母。}\BibleRef{弟前 5:3} 由此我们观察到，寡妇应具备各般亲情：爱自己的子女，并对父母尽其本分。因此，当她向父母尽本分时，也是在教导自己的子女，并且她自己因履行本分而得赏报，因为她为别人所做的，也是为了自己的益处。\par

8. 他说：\BibleQuote{因为这在天主面前是可蒙悦纳的。}\BibleRef{弟前 5:3} 所以，寡妇啊，你若挂虑天主的事，就应当追求你所知道的那些蒙天主喜悦的事。的确，宗徒在前面稍远处\BibleRef{格前 7:34}劝勉寡妇追求\OriginalTerm{节欲}{continence}，说她们挂虑主的事。但在别处，当要拣选一位合格的寡妇时，他吩咐她不仅要记念，更要仰望主：经上说：\BibleQuote{那真正守寡而孤独无依的寡妇，必须仰望天主，并且昼夜不断地恳切祈求与祈祷。}\BibleRef{弟前 5:5} 而他也并非无缘无故地指出，这些人应当无可指摘，因为既然有德性的工作托付给她们，她们也受到了极大的尊敬，甚至为主教所尊重。\par

9. 至于被选定者应是何等样人，导师本人亲口描述了：\BibleQuote{年纪不下六十岁，只作过一个丈夫的妻子。}\BibleRef{弟前 5:9} 这并非说单单年老就足以为寡妇，而是说寡妇的功绩正是老年的本分。因为，那能克制青春的热火、和年轻时代的浮躁激情，不渴望丈夫的温存，亦不贪求子女满堂的丰盛乐趣的寡妇，较比那身体已衰败、年事已高、年岁成熟、既不能再以欢娱取暖，也无望生育后代的寡妇，确然更为高贵。\par

10. 事实上，也没有人被排除在守寡的虔诚之外，如果她在进入再婚之后——对此宗徒的诫命当然没有谴责，好像贞洁的果实就此丧失——然后再次从丈夫的约束中解脱出来。她确实会拥有她的贞洁之功，纵使这功来得迟晚；但那位经历了再婚的寡妇，反而更蒙嘉许，因为在她身上，贞洁的渴望彰明昭著，至于另一类寡妇，则似乎是老年或羞愧使她停止了婚嫁。\par

11. 然而，身体的贞洁并非寡妇的唯一坚定目标，还有更重大、更丰硕的德性操练。\BibleQuote{必须有行善的声望，就如：曾养育儿女，曾款待过旅客，曾为圣人洗脚，曾周济遭难的人，曾尽力从事各种善工。}\BibleRef{弟前 5:10} 你看他包含了多少德性实践！他首先要求虔敬的本分；其次，是款待与谦卑服务的行为；第三，是慈悲与慷慨助人的职务；最后，是力行一切善工。\par

因此，他认为应避免年轻寡妇，\BibleRef{弟前 5:11} 因为她们无法满足如此高超德行的要求。年轻人容易跌倒，因为燃炽的青春热情煽动了种种欲望，而良医的责任在于阻止产生罪的质料。训练灵魂的首要功课在于戒避罪过，其次在于培植德行。然而，宗徒既知\Person{亚纳}——那位八十岁的寡妇——自青年时起便是上主事迹的宣报者，我便不认为他意指将年轻者排除于守寡的虔敬之外，尤其当他说道：\BibleQuote{与其欲火中烧，倒不如结婚。}\BibleRef{格前 7:9} 原来，他确实推崇婚姻作为补救，使那否则将丧亡的人得救；他并未规定能自制者不得度守贞的生活，因为救助跌倒者是一回事，劝人修德行善则是另一回事。\par

关于人的审判，我又能说什么呢？因为在\Person{天主}的审判中，犹太人被指出得罪上主，无非是因为侵犯了寡妇应得的权利和孤儿的权益。先知的呼声宣告，这正是给犹太人招致被遗弃之罚的原因。这也是唯一能平息天主对他们罪恶的义怒的理由——只要他们尊敬寡妇，并为孤儿执行公正的审判。因为经上这样记载：\BibleQuote{你们要维护孤儿，为寡妇辩护，来吧，让我们互相辩论——上主说。}\BibleRef{依 1:17} 另一处又说：\BibleQuote{上主必保护孤儿寡妇。}还有：\BibleQuote{我必丰厚地祝福她的寡妇。} 其中也预表了\OriginalTerm{教会}{Church}的形象。圣洁的寡妇们，你们由此便明白，那因天主恩宠的助佑而受尊崇的职分，绝不可因不洁的欲望而贬损。\par

% structure: p0023 source=h2 target=section reason=relative-heading-rank; base=h2

\section{第三章}

\Person{圣安博罗削}回到\Place{匝尔法特}寡妇的故事，指出她代表\OriginalTerm{教会}{Church}，因而她为贞女、已婚妇女和寡妇的榜样。接着他论及先知预表\Person{基督}，因为先知预示了\OriginalTerm{奥迹}{mysteries}及将临的雨。随后他触及并解释了\Person{基德红}的\Emph{双重标记}，指出行\OriginalTerm{奇迹}{miracles}并非人皆能之；\OriginalTerm{基督的降生成人}{Incarnation of Christ}和犹太人的被弃绝在那叙述中也已预表。\par

回到上文所论，当时遍地大饥荒，然而\Person{天主}对那寡妇的照顾并未缺失，并有\OriginalTerm{先知}{prophet}奉派养活她，这有何意义呢？在这个故事里，\Person{主}提醒我祂要在真理中发言时，\BibleRef{路 4:25} 祂似乎召唤我的耳朵注意一项\OriginalTerm{奥迹}{mystery}。因为有什么比\OriginalTerm{基督}{Christ}与\OriginalTerm{教会}{Church}的奥迹更真实呢？因此，在多数的寡妇中特选一位，并非没有目的。谁是那一位，竟能使一位被提升天的大\OriginalTerm{先知}{prophet}受差引导，尤其是在那时，天闭塞了三年零六个月，遍地大饥荒之际？饥荒遍地，尽管如此，这寡妇却未尝匮乏。这三年是什么？莫非是指\Person{主}来到世上，却在那\OriginalTerm{无花果树}{fig-tree}上找不到果子的那几年吗？正如经上所载：\BibleQuote{看，三年来，我在这无花果树上找果子，竟找不着。}\BibleRef{路 13:7}\par

这无疑就是那寡妇，经上论她说：\BibleQuote{不生育的，喜乐吧！未经过产痛的，欢呼高唱吧！因为被弃者的子女比有夫之妇的子女还多。}\BibleRef{依 54:1} 且她确实是个寡妇，经上又论她如此说：\BibleQuote{你不再记忆你早日的羞辱和你寡居的耻辱，因为我，造你的\Person{上主}，是你的\Person{上主}。}\BibleRef{依 54:4} 或许因此，她是个寡妇，她确实曾因祂身体的苦难而失去丈夫，但在审判之日，必将重新获得那似乎已失去的\OriginalTerm{人子}{Son of Man}。\BibleQuote{我离弃了你只是一会儿，}\BibleRef{依 54:7} 祂说，好使她既被离弃，便更能光荣地坚守\OriginalTerm{信德}{faith}。\par

因此，众人都有可效法的榜样：\OriginalTerm{贞女}{virgins}、\OriginalTerm{已婚妇女}{married women}，以及寡妇。或许\OriginalTerm{教会}{Church}因此既是\OriginalTerm{贞女}{virgin}，又是\OriginalTerm{已婚者}{married}，也是寡妇，因为他们在\Person{基督}内同属一个身体。她正是那位寡妇，为了她的缘故，当地上缺乏\OriginalTerm{天上圣言}{heavenly Word}时，\OriginalTerm{先知}{prophets}们便被设立，因为曾有一位寡妇，并未生育，却将自己的生育保留到适当的时机。\par

因此，那以自己的话用天上的露水滋润干地，且绝非靠人力敞开闭塞之天的人，在我们看来绝非小人物。因为除了\Person{基督}，谁还能敞开天呢？为了\Person{基督}，每天从罪人的食粮中聚集，使\OriginalTerm{教会}{Church}得以增长。因为人没有能力说：\BibleQuote{缸中的面不会用完，罐中的油也不会短缺，直到\Person{上主}使雨降在地上的日子。}即使\OriginalTerm{先知}{prophet}有权这样说话，那声音其实也是\Person{上主}的声音。因此首先声明：\BibleQuote{因为\Person{上主}这样说。}因为确属\Person{上主}的，是保证\OriginalTerm{天上圣事}{heavenly sacraments}的延续，应许\OriginalTerm{属灵喜乐的恩宠}{grace of spiritual joy}不至匮乏，赐予生命的保障、\OriginalTerm{信德}{faith}的印记和\OriginalTerm{德行}{virtues}的恩赐。\par

但\BibleQuote{直到\Person{上主}使雨降在地上的日子}这话是什么意思呢？无非是说祂也要\BibleQuote{像雨落在羊毛上，如降下润泽田地的甘霖}。在这段经文中，揭示了古老历史的\OriginalTerm{奥迹}{mystery}：\Person{基德红}，这位参与奥秘争战的勇士，接受了未来胜利的保证，在心灵的异象中认出了\OriginalTerm{属灵的圣事}{spiritual sacrament}——那雨即是\OriginalTerm{天主圣言}{Divine Word}的露水，它先是落在羊毛上，当时全地正因持续干旱而枯焦；随后又以第二个真实的标记，用阵雨润湿了全地的禾场，而羊毛上却是干的。\par

因为这位有先见之明的人看到了将来教会增长的征兆。因为首先在犹太地区，天主言语的甘露开始滋润（因为\Quote{在犹太，天主为人所认识}），而整个大地却还没有信德的甘露。但当约瑟的羊群开始否认天主，并因冒险犯下种种大罪而在天主前负罪时，那时，当天上雨露的甘露倾注在整个大地上时，犹太人民便开始在自己的不信中干涸枯竭，那时，先知的密云和宗徒们的甘霖浇灌了从世界各地聚集起来的圣教会。这就是那雨水，时而由地上的水分凝聚，时而由山间的雾气凝聚，却借着天上的圣经的救恩甘霖，弥漫整个世界。\par

因此，由这个例子表明，并非所有人都能获得天主大能奇迹的恩宠，只有那些在虔敬的热诚上得到帮助的人才能获得，而那些对上天缺乏敬畏的人，则丧失了天主工作的果实。也以奥迹显示出来：天主子为复兴教会，亲自取了肉身的奥迹，抛弃了犹太民族，从他们那里，谋士、先知和天主恩惠的奇迹都被夺去了，\BibleRef{依 3:2}，因为他们似乎由于某种民族性的污点，不愿意信仰天主子。\par

% structure: p0032 source=h2 target=section reason=relative-heading-rank; base=h2

\section{第4章}

借着亚纳的榜样，圣盎博罗削指明寡妇应有的生活样式，并显示她在各个年龄阶段都是贞洁的典范。由此他论证同一美德有三个等级，都包含在教会内，并举出玛利亚、亚纳和苏撒纳等多种实例。但他补充说，童贞状态优越于其它两种，但寡妇应更加小心维护自己的名声。\par

\Emph{圣经}教导我们，合一带来多大的恩宠，天主对寡妇的祝福又是多么丰厚的恩赐。既然天主如此尊荣她们，我们必须遵守与之相称的生活方式；因为亚纳显示了寡妇应有的样子，她因丈夫早逝而孤苦无依，却获得了圆满赞颂的赏报，她专注于虔敬本分，不亚于追求贞洁。经上说，她守寡八十四年，从不离开圣殿，昼夜以斋戒和祈祷事奉天主。\BibleRef{路 2:36}\par

你可以看到所谓的寡妇是怎样的人：她是一夫之妻，经受了年龄的考验，在虔敬上旺盛，身体却衰残，她的安息之所是圣殿，她的言谈是祈祷，她的生活是斋戒，她日夜以不倦的虔诚事奉，虽然身体承认年老，但在她的虔诚上却不知衰老。这样，寡妇从青年时便已受培养，老年时受到称誉，她保守她的寡居，不是由于时运，也不是因身体软弱，而是借着德性的恢宏。因为经上说她童贞后七年与丈夫一起生活，这正表明，支撑她老年的那些事，在她年轻时的志向中便已开始了。\par

因此我们受到教导，贞洁的美德有三重：一种是婚姻生活的，第二种是寡居的，第三种是童贞的，因为我们不是高举一种而排斥其它。它们各自产生出各自所属的果实。教会的培育在这一点上十分丰富：它拥有可以摆在其他人面前的，却没有可以弃绝的，但愿它绝不会有这样的人！我们讲论童贞，并不是要排斥寡居，我们尊敬寡妇，也为婚姻保留其本身的尊荣。教导这些的，不是我们的规条，而是天主的言语。\par

那么，让我们回想玛利亚、亚纳和苏撒纳是怎样受到颂扬的。既然我们不仅应当颂赞她们，还应效法她们的生活方式，就让我们回想苏撒纳、亚纳\BibleRef{路 2:37}和玛利亚\BibleRef{路 1:28}在哪里被找到，并留意她们各自如何受到特殊的称赞，以及在哪里被提及：那已婚者是在花园里，寡妇在圣殿里，童贞女在密室里。\par

但前者（婚姻）的果实较晚，童贞的果实较早；老年证实她们（寡妇等），童贞却是青年的赞颂，无需年岁的帮助，而是每个年龄的果实。它适宜童年，美化青年，增加老年的尊严，在一切年龄都有它义德的白发、端庄的成熟、端庄的帕子，这并不妨碍虔诚，反而增进虔敬。因为我们从下文看到，圣玛利亚每年在逾越节都与若瑟同往耶路撒冷。\BibleRef{路 2:41} 在童贞女身旁，处处都有热忱的虔诚和贞洁的热心分享者。主的母亲并不因自己的功绩稳妥而自负，她越认识自己的功绩，就越充分地还愿，越丰盛地服务，越圆满地尽职，越虔敬地履行本分并充实奥迹的时期。\par

那么，你们更应当专心追求贞洁，免得你们给那些仅有端庄和品行作证的人留下任何非议的余地。因为一位童贞女，虽然在她身上，品格比身体更为重要，但她以身体的完整消除了诽谤；而一位寡妇，既已失去能证明童贞的帮助，她接受贞洁的审断，不是凭接生婆的话，而是凭她自己的生活方式。因此，圣经已指明，寡妇的性情该是多么谨慎和虔敬。\par

% structure: p0040 source=h2 target=section reason=relative-heading-rank; base=h2

\section{第5章}

福音中那位寡妇的榜样推荐了济贫的慷慨，她两个小钱被看重，胜过富人的大量奉献。这两个小钱被神秘地解释为代表着两约。我们被教导要奉献的宝藏是什么？按照贤士们的榜样，要献上三份礼物；或按照寡妇的榜样，献上两份。圣盎博罗削在本章结尾劝勉寡妇要热心于善工。\par

27. 在同一卷书的另一处，我们也受到教导：向穷人施行仁慈与慷慨是多么合宜，而这种心肠不应顾忌自身的贫穷而受遏制，因为慷慨并不取决于财物的多寡，而在于施舍的心意。因为按主的声音，那个寡妇被置于众人之上，论到她曾说：\BibleQuote{这个寡妇投的比众人都多。}\BibleRef{路 21:3} 主在此事例中特别教导众人，谁也不可因羞愧自己的贫穷而退缩不行施舍，富人也不可自诩其捐助似乎多于穷人。因为从小额积蓄中拿出的一个小钱，比富余中取出的财宝更丰富，因为所考虑的不是所给的数额，而是所剩余的数额。没有人比那不留给自己什么的人施舍的更多。\par

28. 富有的妇人啊，你为何与穷人相比而夸耀自己？当你周身披戴黄金，拖着贵重的长袍，却渴望受人尊敬，仿佛她因与你的财富相比而显得卑微渺小，因为你用你的施舍胜过了穷乏人？江河盈满时也会泛滥，但小溪的涓流更为甘甜。新酒发酵时泡沫翻腾，农夫并不认为溢出的是损失。打谷时，籽粒从呻吟的禾场上掉落；然而即使庄稼歉收，面缸中的面粉并不减少，满瓶的油倾流而出。但饮宴却倾空了富人的酒桶，而寡妇的小油瓶却涌出丰盛。那么，你所当算为奉献的，是你出于虔敬所给予的，而不是你轻蔑地掷出的。总之，没有人比那用儿女的食粮供养先知的她施舍的更多。因此，既然无人施舍得更多，也就无人有更大的功德。这具有道德上的训诲意义。\par

29. 若论奥秘的意义，我们不可轻看这位往\OriginalTerm{银库}{treasury}里投了两个\OriginalTerm{小钱}{two mites}的妇人。显然，这位妇人是高贵的，在天主的判断中堪受优于众人的称许。或许她就是那凭信德献上\OriginalTerm{两约}{testaments}以帮助世人的那位，因此无人比她的施舍更大。也没有人能比得上她的捐献之量，因为她将信德与仁慈结合在一起。那么，无论你是谁，在生命中实践寡居生活的，不要迟疑，将充满信德与恩宠的两个小钱投进银库。\par

30. 那从自己的宝藏中取出君王的完美肖像的妇人是有福的。你的宝藏是智慧，你的宝藏是贞洁与正义，你的宝藏是良好的明悟，正如那宝藏，当\Person{贤士}们朝拜主时，从中取出黄金、\OriginalTerm{乳香}{frankincense}和\OriginalTerm{没药}{myrrh}；\BibleRef{玛 2:11} 用黄金标示君王的权能，用乳香敬拜天主，用没药宣认身体的复活。你们若省察自己，也有这宝藏：\BibleQuote{因为我们是在瓦器中存有这宝贝。}\BibleRef{格后 4:7} 你们有可施舍的黄金，因为天主向你们索取的，不是闪闪发光金属的贵重礼物，而是那在审判之日，烈火不能烧毁的黄金。祂也不求珍贵的礼物，而是信德的芬芳，这芬芳从你们心之祭台上发出，由虔诚心灵的气息呼出。\par

31. 那么，从这宝藏中，不仅贤士的三件礼物，寡妇的两个\OriginalTerm{小钱}{two mites}也得以取出，在这小钱上，闪耀着天上君王的完美肖像，祂光荣的光辉和祂本体的肖像。那藉着劳累和日常劳作，寡妇以贞洁换来的艰辛收获也是珍贵的，她日夜不断地操劳，以有益贞洁的警醒劳作积蓄财宝，为能保全亡夫的床榻不受玷污，能支持她可爱的儿女，并服事穷人。她优先于富人，她将不畏惧基督的审判。\par

32. 我的女儿们，努力效法她吧：\BibleQuote{对善事热心原是好的。}\BibleRef{迦 4:18} \BibleQuote{你们该热切追求那更大的恩赐。}\BibleRef{格前 12:31} 主时常注视着你，耶稣走向银库时注视着你，你也想着从你善工的所得中要给予急需者帮助。那么，这又算什么呢，你献上你的两个小钱而获得\OriginalTerm{主的身体}{Lord's Body}？所以，不可空手到主你的天主面前去，\BibleRef{出 34:20} 不可缺乏仁慈、不可缺乏信德、不可缺乏贞洁；因为主耶稣惯常注视并称赞的，不是空虚的人，而是富有德行的。让少女看见你在工作，让她看见你在服事他人。因为这是你向天主应尽的回报，你应当从他人的进步中向天主献上回报。没有比虔诚的奉献更蒙天主悦纳的回报了。\par

% structure: p0048 source=h2 target=section reason=relative-heading-rank; base=h2

\section{第6章}

\Person{纳敖米}是一个寡妇从儿媳那里收获自己良好训练之果实的实例，也是一个标记，表明善良的寡妇绝不会缺乏必要的供养。即使她的生活看似凄凉，她却是幸福的，因为主的许诺是给她的。\Person{圣盎博罗削}进而提及哭泣的益处。\par

33. \Person{纳敖米}寡妇在你看来不算什么吗？她靠拾取他人田间的麦穗维持寡居，在垂老之年由儿媳供养？\BibleRef{卢 2:2} 对于寡妇的供养与益处而言，她们如此训练儿媳，以致在垂暮之年能从她们那里得到扶持，并作为教导的酬报和训练的赏报，这是极大的神益。因为对于一个善于教导、好好训诲儿媳的人来说，绝不会有缺一个\Person{卢德}，她宁愿选择婆婆的寡居生活，胜过自己的父家，即使丈夫已死，也不离开她，在急需时供养她，在忧苦中安慰她，即便被遣走也不舍弃她；因为良好的教训永远不会匮乏。所以，纳敖米失去了丈夫和两个儿子，失去了她繁殖的果实，却没有失去她虔诚关怀的赏报，因为她既在忧苦中找到了安慰，也在贫困中得到了扶持。\par

34. 圣洁的妇女们，那么你们看，寡妇在美德的子女和她的功劳的果实上是多么丰饶，这些果实永无止境。因此，一位好寡妇一无所缺；即使她因年老而疲惫，处于极度的贫困中，她通常也会得到她所给予教导的报酬。即使最亲近的人辜负了她，她也会找到那些不那么亲近的人来孝敬她们的母辈、尊敬她们的长辈，并通过微薄的赡养礼物渴望获得他们自己善行的果实，因为给予寡妇的礼物必得丰厚的回报。她求食，却回报珍宝。\par

35. 但她似乎度着悲伤的日子，在泪水中度过时光。她由此更为有福，因为她藉着少许哭泣为自己买得永恒的喜乐，以片刻的代价获得永生。对于这样的人，经上说得好：\BibleQuote{哀恸的人是有福的，因为他们要欢笑。}\BibleRef{路 6:21} 那么，谁会宁愿选择现今欢乐的虚假外表，而不取将来无忧无虑的愉快呢？按肉身是主的蒙选先祖，他吃灰如吃饼，以眼泪调和他的饮料，并藉着夜间的眼泪在清晨为自己赢得了救赎的喜乐，难道他在我们看来是微不足道的权威吗？他从哪里获得那极大的喜乐，岂不是因为他大大哭泣，并以自己的眼泪为代价，为自己获得了未来光荣的恩宠吗？\par

36. 那么，寡妇有这美好的推荐：她在哀悼丈夫的同时，也为世界哭泣，那救赎性的眼泪已经预备好，为逝者流下，必有益于生者。眼中的泪水与心中的忧伤相合，它激起怜悯，减轻劳苦，缓解悲伤，并保持端庄；她不再觉得自己如此可怜，因为在眼泪中找到了安慰，这些眼泪是爱的报偿，也是虔诚记忆的证明。\par

% structure: p0054 source=h2 target=section reason=relative-heading-rank; base=h2

\section{第七章}

藉着\Person{友弟德}的榜样表明，寡妇并不缺乏勇气；她拜访\Person{敖罗斐乃}之前的准备被详细叙述，同样也提及了她的贞洁、智慧、节酒和节制。最后，\Person{圣盎博罗削}在证明她的勇敢不亚于她的明智之后，陈述了她成功之后的端庄。\par

37. 但一位好寡妇通常也不缺乏勇敢。因为真正的勇敢在于，藉着心灵的虔诚，超越通常的天性和女性的软弱，正如那位名叫\Person{友弟德}的女子所拥有的。她仅凭自己一人，就能唤醒那些因围困而崩溃、因恐惧而惊骇、因饥饿而憔悴的人们，并保护他们免受敌人的侵害。因为，如我们读到的，当\Person{敖罗斐乃}在许多战役中得胜后令人畏惧，把无数人驱入城墙之内；当武装的男子们感到害怕，并且已经在商讨最后的投降时，她走出城墙，既超越了她所解救的那支军队，又比她那支驱散的军队更为勇敢。\par

38. 但为了了解成熟寡妇身份的特质，请通读\WorkTitle{圣经}的历程。自她丈夫去世以后，她便脱去喜乐的衣裳，穿上哀悼的服饰。除了安息日、主日和节庆日之外，她每天都专务守斋，这不是因为屈从于恢复体力的欲望，而是出于对宗教的尊重。因为经上这样说：\BibleQuote{你们或吃或喝，无论做什么，都要因主耶稣基督的名而行。}\BibleRef{格前 10:31} 意思是，就连身体的休养，也应顾及神圣宗教的敬拜。因此，\Person{圣友弟德}因长期的哀悼和每日的守斋而坚强，不寻求世界的享乐，不顾危险，并对死亡无所畏惧。为了实施她的计谋，她穿上那件喜乐的外袍，就是她丈夫在世时她常穿的那件，好像她若拯救了国家，便能使丈夫喜悦一样。但她看到了另一位男子，就是她想要取悦的那一位，也就是经上论及的那一位：\BibleQuote{在我以后来的那位，成了在我以前的。}\BibleRef{若 1:30} 她重戴婚嫁的装饰去作战，做得很好，因为婚姻的提醒是贞洁的武器，除此之外，寡妇无法取悦他人或获得胜利。\par

39. 又何必叙述后续呢？她在成千上万的敌人中保持了贞洁。又何必讲论她的智慧，说她如何设计了这样的计谋呢？她选择了那位统帅，以免受下属的侮辱，并为胜利预备机会。她保留了节食的功劳和贞洁的恩宠。因为，如我们所读到的，她既未因食物也未因奸淫而被玷污，她藉着保持贞洁而战胜敌人，丝毫不亚于她拯救国家所获得的胜利。\par

40. 对于她的节酒，我还有什么可说的呢？节德确实是女性的美德。当男人们被酒灌醉、沉睡不醒时，寡妇拿起剑，伸手砍下战士的头，毫发无伤地穿过敌人的阵中。那么，你们要注意，醉酒能给女人带来多大的伤害，因为酒如此削弱男人，以致他们被女人战胜。因此，寡妇应当节制，首先要戒酒，好能戒绝奸淫。若酒不诱惑你，他诱惑你也徒劳。因为如果\Person{友弟德}喝了酒，她就会与奸夫同睡。但正因她不喝酒，一个人的节酒便毫无困难地能够战胜并逃离醉酒的军队。\par

41. 而这更多不是她双手的功劳，反而更是她智慧的奖杯。因为她仅凭一己之力战胜了\Person{敖罗斐乃}，又以自己的智慧战胜了敌人的全军。因为悬挂\Person{敖罗斐乃}的头颅，这一壮举是男人的智慧未能策划的，她鼓舞了同胞的勇气，瓦解了敌人的士气。她以自己的端庄激励了自己的朋友，并使敌人惊恐，以致他们溃逃并被杀。因此，一位寡妇的节德和节酒不仅克制了她的本性，而且，更甚的是，甚至使男人们更加勇敢。\par

42. 然而她并没有因这次成功而得意忘形，尽管她本可因胜利而欢欣鼓舞，她却没有放弃守寡的操练，反而拒绝了所有想娶她的人，放下了欢乐的衣裳，重新穿上了寡妇的服装，不在意胜利的装饰，认为克制肉身的恶习比战胜敌人的武器更好。\par

% structure: p0062 source=h2 target=section reason=relative-heading-rank; base=h2

\section{第八章}

虽然许多其他寡妇在德行上接近\Person{友弟德}，圣\Person{盎博罗削}却只提议谈论\Person{德波辣}。她必定是寡妇们的一个德行楷模，竟被拣选去治理和保卫男子。她的一大荣耀是，当她儿子统率军队时，他拒绝出征，除非她也同去。因此她领导军队并预言了结果。在这个故事中，展示了教会的冲突与胜利，以及她的属神的武器，并从妇女身上除去了软弱的一切借口。\par

43. 为了不让人觉得似乎只有一个寡妇完成了这无可比拟的工作，毫无疑问还有许多其他德行相等或几乎相等的妇女，因为良种通常结出许多充满谷粒的穗子。所以，不要怀疑那古老的耕种期在众多妇女的品格上是丰收的。但若要囊括所有人则过于冗长，就细想几个，尤其是\Person{德波辣}，圣经为我们记述了她的德行。\par

44. 因为她不仅表明了寡妇无需男人的帮助，因为——毫不受性别软弱的约束——她承担了男人的职责，而且所做超乎所承。最后，当犹太人在民长领导下被治理时，因为他们不能用男子汉的正义治理他们，也不能用男子汉的力量保卫他们，以致四方战火燃起，他们便拣选了\Person{德波辣}，好按她的判断受治理。这样，一个寡妇既在和平中治理成千上万的男子，又保卫他们脱离仇敌。在\Place{以色列}有许多民长，但在她之前没有一个女子作民长，就如\Person{若苏厄}之后有许多民长，却没有一个是先知。我以为，她的民长职位被记述，她的事迹被描述，就是为了让妇女不要因性别的软弱而被拦阻于英勇的作为。一个寡妇，她治理百姓；一个寡妇，她统率军队；一个寡妇，她拣选将领；一个寡妇，她决定战争并安排凯旋。所以，当受指责的或易于软弱的，并非天性。使人坚强的，不是性别，而是英勇。\par

45. 在和平时期，对这妇人毫无怨言，也找不到什么过失，而大多数民长却给了百姓不小的罪过。但当客纳罕人——一个战斗凶猛、兵多将广的民族——相继联合起来，对犹太人显出可怕的样子时，这位寡妇先于所有人，作好了一切战争准备。为表明家庭所需并不依赖公共资源，反而公共职责受家庭生活纪律的引导，她从家里领出她的儿子作军队的统帅，好让我们承认一位寡妇能训练战士；她身为母亲教导他，身为民长任命他统率，由于自身勇猛，她训练他，身为女先知，她差遣他去获取必定胜利。\par

46. 最后，她的儿子\Person{巴辣克}表明了胜利的主要部分在于妇人手中，当他说：\BibleQuote{你若不同我去，我也不去；因为我不知主派他的使者与我同去的那一天。} 那么，这位妇人是何等大能，军队的统帅竟对她说：\BibleQuote{你若不同我去，我也不去。} 我说，这位寡妇的刚毅何等大，她不因母爱而让儿子远离危险，反而以母亲的热忱劝勉儿子去争取胜利，同时说那场胜利的决定权在于一位妇人手中！\par

47. 因此，\Person{德波辣}预言了战事的结果。\Person{巴辣克}，依命率领军队出战；\Person{雅厄耳}取得了凯旋，因为\Person{德波辣}的预言为她作战，这预言在奥迹中向我们揭示了教会从外邦人中兴起，为她——教会——找到了战胜\Person{息色辣}，即敌对的势力的凯旋。因此，是先知们的谕言为我们作战，是先知们的那些判断和武器为我们赢得了胜利。为此，战胜仇敌的，不是犹太人，而是\Person{雅厄耳}。所以，那个民族是不幸的，它不能凭信德去追击它所击溃的敌人。因此，由于他们的过错，救恩临到了外邦人，由于他们的迟钝，胜利留给了我们。\par

48. 于是\Person{雅厄耳}消灭了\Person{息色辣}，然而，\Person{息色辣}是被犹太老兵队伍在其杰出领袖率领下击溃的，因为\Person{巴辣克}这名字有如此含义；如我们常读到的那样，先知们的言论和功绩为圣祖们获取了天上的援助。但即使在那时，对抗\OriginalTerm{邪恶势力}{spiritual wickedness}的胜利已为那些在福音中听到说：\BibleQuote{我父所祝福的，你们来承受自创世以来为你们预备的国度吧。}\BibleRef{玛25:34} 的人预备。这样，胜利的开端来自圣祖们，其终局却在教会。\par

49. 但教会制胜仇敌的势力，不是靠这世界的武器，而是靠属神的武器，\BibleQuote{就是凭天主具有大能，足以攻破坚固的堡垒，击破\OriginalTerm{邪恶势力}{spiritual wickedness}的高垒。}\BibleRef{格后10:4} 而\Person{息色辣}的干渴被一碗奶解除了，因为他被智慧所胜，因为对我们作食物有益的东西，对仇敌的势力却是致命和软弱的。教会的武器是信德，教会的武器是祈祷，这祈祷克胜仇敌。\par

50. 所以，依照这段历史，一个妇人——为激发妇女们的心志——成了民长，一个妇人整顿了一切，一个妇人说了预言，一个妇人凯旋了，她加入战阵，教导男子们在一个妇人的领导下作战。但在奥迹中，这是信德的战斗和教会的胜利。\par

51. 所以啊，你们这些妇女，不能以天性作借口。你们这些寡妇，不能以女性软弱作借口，也不能将你们的易变归咎于失去丈夫的扶持。只要灵魂不缺乏勇毅，人人都有足够的防护。对寡妇而言，年老本身便是守贞的普遍保障；对已故丈夫的哀悼、日常劳作、料理家务、为子女操心，常能阻挡有害灵魂的淫乱；而丧服本身、葬礼的肃穆、不断的哭泣、深深刻在悲伤额头上的哀愁，能抑制淫荡的目光，扼止贪欲，使轻浮的眼神转去。充满悔恨的深情，是贞洁的优良守卫；若保持警觉，罪过便无从侵入。\par

% structure: p0073 source=h2 target=section reason=relative-heading-rank; base=h2

\section{第九章}

对于一种反对意见，认为寡妇若能处于愉快的环境中，则守寡或许还可以忍受，\Person{圣盎博罗削}回答道，愉快的境遇比困苦更为危险；并引用\WorkTitle{圣经}中的例子指出，寡妇可以从子女和女婿身上获得许多幸福。她们应当求助于能够帮助我们的宗徒们，并应恳求天使和殉道者的代祷。然后，他论及关于孤独以及照顾财产的一些怨言，最后指出，一个寡妇若已有已婚或适婚年龄的女儿，再婚是不合宜的。\par

52. 所以，你们这些寡妇，已经明了你们并不缺乏自然的帮助，并且能够保持健全的劝导。再者，你们在家中并非没有保护，因为你们甚至能要求公共权力的最高点作为帮助。\par

53. 但也许有人会说，寡妇若生活富裕，守寡还较容易忍受；至于穷困潦倒的寡妇，很快就会被击垮，轻易屈服。关于这一点，我们不单从经验得知，享乐对寡妇而言，比困苦更危险；而且\WorkTitle{圣经}中的例子也告诉我们，即使在软弱中，寡妇通常也不会得不到援助\BibleRef{弟前 5:16}，并且，如果她们很好地抚养了子女，选择了合适的女婿，天主的助佑和人的支援便比其他人更易临于她们。最后，当\Person{西满}的岳母身患重热症卧病时，\Person{伯多禄}和\Person{安德肋}为她求主：\BibleQuote{他遂上前去，站在她身边，叱退热症，热症就离开了她；她立刻起来服事他们。}\BibleRef{路 4:39}\par

54. 经上说：\BibleQuote{她正患着严重的热症，他们为她求他。}\BibleRef{路 4:38}你们也有接近的人为你们祈求。你们有宗徒们接近，你们有殉道者接近；若你们在虔敬上与殉道者结合，也要藉着慈善行为接近他们。你们施行仁慈，就会靠近\Person{伯多禄}。使人接近的，不是血缘关系，而是德行的契合，因为我们行事不是随从肉身，而是随从圣神。所以，要珍惜与\Person{伯多禄}的亲近，与\Person{安德肋}的契合，好让他们为你们祈祷，使你们的贪欲退去。你们，虽躺在世上，一旦被天主的话触动，便会立即起来服事基督。\BibleQuote{因为我们的家乡原是在天上，我们从那里也等候救主，主耶稣基督。}\BibleRef{斐 3:20}因为没有人能躺着服事基督。你们服事穷人，就是服事了基督。他说：\BibleQuote{凡你们对我这些最小兄弟中的一个所做的，就是对我做的。}\BibleRef{玛 25:40}寡妇们，你们若为自己选择这样的女婿，为后代选择这样的恩主和朋友，便有援助。\par

55. 所以，\Person{伯多禄}和\Person{安德肋}为那位寡妇祈祷。但愿有人能如此迅速为我们祈祷，或者更好的是，为岳母祈祷的\Person{伯多禄}和他的兄弟\Person{安德肋}本人。那时他们能为亲戚祈祷，如今他们能为我们和众人祈祷。因为你们见到，被大罪束缚的人，不适合为自己祈求，当然更不易为自己求得。所以她应利用别人替她向那位医生祈求。因为患病的人，除非藉他人的祈祷请来医生，便不能为自己求。肉身是软弱的，灵魂染病，被罪恶的锁链束缚，无法把软弱的脚步引向那位医生的宝座。必须为我们恳求天使，他们曾是我们的护卫；必须恳求殉道者，我们似乎藉着他们身体的遗骸作为抵押，而要求他们的护佑。他们能为我们的罪过祈求，因为他们若有任何罪过，已用自己的血洗净了；因为他们是天主的殉道者，我们的前驱，我们生活和行动的鉴察者。我们不应以请他们作我们软弱的代祷者为耻，因为他们自己也曾体认肉身的软弱，即便他们已得胜了。\par

56. 所以，\Person{伯多禄}的岳母找到了为她祈祷的人。而你们，寡妇啊，如果如同真正孤独无依的寡妇那样仰望天主，不断祈求，恒心祈祷\BibleRef{弟前 5:5}，你们也要找到为你们祈祷的人；每天训练你们的身体如同已死，好使你们死去能再活过来；逃避享乐，这样，你们虽患病，也能得到痊愈。\BibleQuote{但那放纵享乐的，虽生犹死。}\BibleRef{弟前 5:6}\par

57. 你们不再有任何结婚的理由了，你们有人为你们转求。不要说：\Quote{我孤苦伶仃。}这是想结婚者的怨言。不要说：\Quote{我孤单。}贞洁寻求独处：端庄者寻求清静，不端庄者寻求结伴。但你们有必要的俗务；你们也有一位为你们辩护的。你们害怕对头；主将亲自干预判官，说：\BibleQuote{应为孤儿伸冤，为寡妇辩护。}\BibleRef{依 1:17}\par

58. 但你愿意看顾你的产业。端庄的产业更加伟大，而寡妇比已婚者更能守护它。一个奴隶做错了事。宽恕他吧，因为你容忍别人的过失，总比揭露它更好。但你想要结婚。就这样吧。单纯的欲望并不是罪。我不问原因，为何要编造一个呢？你若觉得好，就说出来；若不合适，就保持沉默。不要责怪天主，不要责怪你的亲属，说没有人保护你。但愿这愿望不会落空！也不要说你是在为子女的利益着想，你正在使他们失去母亲。\par

59. 有些事在抽象上是允许的，但因年龄缘故却不可行。为何母亲的婚礼与女儿的婚礼同时筹备，甚至往往更晚？为何成年的女儿学会在母亲的未婚夫面前，而不是在她自己的未婚夫面前脸红？我承认我曾劝你更换衣裳，但不是要你披上新娘的面纱；是要你远离坟墓，而不是准备婚床。一个已有女婿的新婚妇人，这是什么意思？有比孙辈还年幼的子女，是多么不合体统！\par

% structure: p0083 source=h2 target=section reason=relative-heading-rank; base=h2

\section{第十章}

\Person{圣安博罗削}再次回到基督的主题，谈论祂在一切苦难中的美善。那位良医治疗我们疾病的各种方式，以及只要我们不忘呼求祂，医治便会迅速。他触及意志的道德意义，并指出这在\Person{伯多禄}的岳母身上彰明显现；最后指出基督的仆役、尤其是主教当如何行事，并说他们尤其必须靠着恩宠兴起。\par

60. 但让我们回到正题，不要在自己罪过的伤痛中离开医生，在照顾他人的疮伤时，让自己的疮伤继续加重。那时，人们求问这位医生。不要因为主是伟大的，就害怕祂或许不肯屈尊来看顾病人，因为祂常从天上来到我们这里；祂惯常探望的不仅是富人，还有穷人和穷人的仆人。\BibleRef{路 4:18} 因此，当有人呼求时，祂现在也临到\Person{伯多禄}的岳母。\BibleQuote{祂站在她身旁，斥责了热症，热症就离了她；她立刻起身服事他们。}\BibleRef{路 4:38} 祂既配得纪念，也配得渴慕，配得爱戴，因为祂眷顾与人相关的每一件事，祂的作为奇妙。祂不嫌弃探望寡妇，进入穷苦人家狭窄的房间。作为天主，祂施命；作为人，祂探望。\par

61. 感谢\WorkTitle{福音}，借着它，我们这些在基督来到这世界时未曾看见祂的人，在阅读祂的事迹时，就好像与祂同在；正如那些祂走近的人从祂那里获取了信德，愿祂在我们相信祂的事迹时，也走近我们。\par

62. 你看见与祂同在的是怎样的医治吗？祂斥责热症，祂命令邪魔，在别处祂按手在他们身上。那时，祂不仅用言语，也藉触摸治愈病人。那么你呢，你若为许多欲望所焚烧，或因某人的美貌或财富而被掳获，就应当恳求基督，呼求医生，向祂伸出你的右手，让天主的手触摸你的内心深处，让天上圣言的恩宠进入你内心欲望的脉络，让天主的右手触摸你心灵的秘密。祂在一些人的眼上敷了泥，好叫他们看见，\BibleRef{若 9:6} 万物的创造者教训我们：应当记念自己的本性，认清我们肉身的卑贱；因为除非人因认识自己的卑贱而不自高自大，没有人能看见神圣的事物。另一个人被吩咐去给司铎查看，好叫他永远摆脱癞病的鳞屑。\BibleRef{路 5:14} 唯独这样的人才能力保自己身体和灵魂的纯洁：他知道怎样将自己呈明于那位司铎，就是我们所得以为我们罪过的中保，并清楚地对祂说：\BibleQuote{你按照\Person{默基瑟德}的品位，永为司铎。}\par

63. 不要害怕医治会有任何迟延。被基督治愈的人毫无障碍。你必须使用已获得的药方；祂一发出命令，盲人便看见，瘸子行走，哑巴说话，聋子听见，患热症的妇人起来服事，附魔者获得释放。那么你呢，你若有任何不体面的欲望，为渴望某事而憔悴，就当恳求主，向祂显示你的信德，不要怕迟延。哪里有祈祷，圣言就在那里，欲望便被逐出，邪情便离去了。不要害怕因告解而触犯，反要视之为权利，因为你这先前因身体的重病而受苦的人，将要开始服事基督。\par

64. 在这里还可以看到\Person{伯多禄}岳母的意志倾向，她由此为自己领受了，仿佛未来之事的种粒，因为对每个人来说，意志是他未来之事的原因。因为智慧从意志涌出，智慧人娶她为妻，说：\BibleQuote{我切愿聘娶她。}\BibleRef{智 8:2} 那么，这个意志起初在各种欲望的热症下软弱无力，后来因着宗徒们的职务而坚强起来，去服事基督。\par

65. 同时，这也显示出服事基督者应当是怎样的人，因为他首先必须摆脱各种快乐的引诱，必须摆脱身体和灵魂内在的软弱，才能去分施基督的圣体圣血。因为凡因自己的罪而患病、远非健全的人，不能分施永生之医治的药方。司铎啊，你要看清你所行的事，不可用发烧的手触摸基督的圣体。你当先得医治，才能去服事。如果基督吩咐那些如今洁净了、但曾患过癞病的人去给司铎们查看，\BibleRef{路 17:14} 那么司铎本人保持洁净，岂不更应当吗？因此，那位寡妇不能因我没有顾惜她而生气，因为我也没有顾惜自己。\par

\Person{伯多禄}的岳母，如经上记载的，起来服侍他们。说\BibleQuote{起来}确实合宜，因为\OriginalTerm{宗徒职分}{apostleship}的恩宠已预示着圣事的预像。基督的仆人理应起来，正如经上所写：\BibleQuote{你这睡着的人，应当醒过来，从死者中起来！} \BibleRef{弗 5:14}\par

% structure: p0092 source=h2 target=section reason=relative-heading-rank; base=h2

\section{第11章}

已经表明通常为再婚而提出的借口没有分量，\Person{圣安博罗削}宣称他不谴责再婚，尽管从宗徒的话中他指出了其不便之处，虽然那些再婚者在教会中是允许的，他趁机提及禁止再婚的\OriginalTerm{异端者}{heretics}。并且他说，因为不同人的力量各异，所以\OriginalTerm{洁德}{chastity}不是命令，只是推荐。\par

因此我说，那些习惯施予的寡妇，既不缺乏必要的开销，也不缺乏帮助；她们在极大的危险中常护守丈夫的资财；再者，我认为丈夫的善举通常由女婿和其他亲属弥补给她们，而天主的怜悯更乐于扶助她们；因此，若无特殊原因再醮，便不应存有这等愿望。\par

然而，我这话是作为\OriginalTerm{劝告}{counsel}，并不是以\OriginalTerm{命令}{precept}吩咐，为激发寡妇们的意愿，而非约束她们。因为我不禁止再婚，只是不建议。人性软弱的考量是一回事，\OriginalTerm{洁德}{chastity}的恩宠是另一回事。我更进一步说，我不禁止第二次婚姻，但不赞成屡次结婚，因为并非所有合法的事都有益：宗徒说：\BibleQuote{凡事我都可行，} \BibleQuote{但不都有益。} \BibleRef{格前 6:12} 又如，饮酒是合法的，但往往并不有益。\par

因此，结婚是合法的，但戒绝更为合宜，因为婚姻中有\OriginalTerm{束缚}{bond}。你要问是怎样的束缚？\BibleQuote{凡有丈夫的妇人，丈夫活着的时候，是被法律束缚的；倘若丈夫死了，她便脱离了丈夫的法律。} \BibleRef{罗 7:2} 这就证明婚姻是束缚，妇人为其所缚，也得脱离。相互爱情的恩宠虽美，但\OriginalTerm{奴役}{servitude}更为恒常。\BibleQuote{妻子对自己的身体没有主权，而是丈夫有主权。} \BibleRef{格前 7:4} 为使这奴役不似出于性别，而非婚姻，接着说：\BibleQuote{同样，丈夫对自己的身体也没有主权，而是妻子有主权。} 那么，婚姻的约束何其大，竟使强者也屈服对方；因相互约束，各人必须服事。若一方愿意节制，也不能从轭下缩颈避开，因为他受制于对方的无节制。经上说：\BibleQuote{你们是用高价买来的，切不要做人的奴隶。} \BibleRef{格前 7:23} 你看婚姻的奴役界定得何等明白。这不是我说的，是宗徒说的；更确切地说，不是他，而是在他内说话的基督。他论及善度婚姻者的这种奴役。因为前面你们读到：\BibleQuote{不信主的丈夫因信主的妻子成了圣洁的；不信主的妻子也因信主的丈夫成了圣洁的。} \BibleRef{格前 7:14} 再往下：\BibleQuote{但若不信主的一方要离去，就由他离去；在这种情形下，兄弟或姊妹不必受拘束。} \BibleRef{格前 7:15} 那么，若善的婚姻是奴役，恶的婚姻又当如何？他们既不能彼此圣化，反而彼此毁灭。\par

然而，我劝勉寡妇们保守所受恩赐的\OriginalTerm{恩宠}{grace}，我也同样激励妇女们遵守教会的纪律，因为教会由众人组成。基督的羊群中，有些需食硬粮，有些仍需哺以乳汁；她们必须提防那些披着羊皮的狼，他们外表装出全然\OriginalTerm{节德}{continence}的模样，却煽动污秽的无节制。他们知道\OriginalTerm{洁德}{chastity}的担子何其沉重，因为他们连指尖也不能触及；他们要求别人超越限度，自己却连任何限度也无法遵守，反而在残酷的重压下屈服。因为担子的份量总应按负者的力量而定；否则，负者软弱，必被加于其上的重担压垮；正如硬食会噎住婴儿的喉咙。\par

因此，正如众多负者中，其力量不凭少数人估量；强者也不会按他人的软弱领受任务，而是各人可按己意承担愿负的重担，酬报随力量增长而加增；同样，不可为妇女设下圈套，也不可将超过其力量的\OriginalTerm{节德}{continence}重担强加于她们，而应让各人自己衡量此事，不是受命令的权威强迫，而是因\OriginalTerm{恩宠}{grace}的增长激励。因此，按德行的不同程度，设定的酬报也不同，一件事不被贬低，另一件事方可受赞扬；但一切都已说明，为使人选那更好的。\par

% structure: p0099 source=h2 target=section reason=relative-heading-rank; base=h2

\section{第12章}

此处论述\OriginalTerm{命令}{precept}与\OriginalTerm{劝告}{counsel}之间的区别，如福音中少年人的事例所示，并论及为\OriginalTerm{劝告}{counsel}和\OriginalTerm{命令}{precept}所设定的不同酬报。\par

72. 因此，\OriginalTerm{婚姻}{marriage}是可敬的，但\OriginalTerm{贞洁}{chastity}更为可敬，因为\BibleQuote{那将自己的童贞女嫁人的，做得好；那不将她嫁人的，做得更好。}\BibleRef{格前 7:28} 那么，好的不必回避，但更好的则应选择。因此，这并不是强加于任何人的，而是摆在他面前的。所以，宗徒说得好：\BibleQuote{关于童贞女，我没有主的命令，但我提出我的~\OriginalTerm{劝谕}{counsel}。}\BibleRef{格前 7:25} 因为命令是向属下发出的，劝谕则是对朋友提出的。哪里有命令，那里就有\OriginalTerm{法律}{law}；哪里有劝谕，那里就有\OriginalTerm{恩宠}{grace}。命令是为强制人遵守合乎本性的事，劝谕是为激励我们追随恩宠。因此，法律是赐给犹太人的，但恩宠则保留给\OriginalTerm{选民}{elect}。法律被赐下，是为藉着对惩罚的恐惧，将那些偏离本性界限的人召回遵守；而恩宠则是为激励选民，既藉着对美善的渴望，也藉着所应许的赏报。\par

73. 如果你记得福音中那个人的事例，就会明白\OriginalTerm{诫命}{precept}与\OriginalTerm{劝谕}{counsel}的区别。那人先被命令不可杀人，不可奸淫，不可作假见证；因为那是带有惩罚的诫命，违反必受罚。但当他说自己遵守了法律的一切诫命后，主又给了他一个劝谕，要他变卖所有的一切，跟随主，\BibleRef{玛 19:18} 因为这些事不是作为命令强加的，而是作为劝谕提出的。因为命令有两种方式，一种是通过诫命，另一种是通过劝谕。因此，主一方面说：\BibleQuote{你不可杀人，}这是给出诫命；另一方面又说：\BibleQuote{你若愿意是成全的，去变卖你所有的。}那么，凡留有选择余地的人，就不受诫命的约束。\par

74. 因此，那些遵行了诫命的人可以说：\BibleQuote{我们是无用的仆人，我们不过做了我们应做的事。}\BibleRef{路 17:10} 但童贞女不会这样说，那变卖了自己一切财物的人也不会这样说，他们反而期待那储存的赏报，就像那位圣宗徒所说的：\BibleQuote{看，我们舍弃了一切跟随了你，那么，将来我们可得到什么呢？}\BibleRef{玛 19:27} 他并不像那无用的仆人，说自己做了应做的事，而是作为对主人有益的，因为他通过所获得的盈利，使托付给他的~\OriginalTerm{塔冷通}{talent}增加了，凭着无愧的良心，毫不焦虑自己的功劳，期待着信德与德行的赏报。因此，主对他说，也对其他人说：\BibleQuote{你们这些跟随我的人，在重生的时代，当人子坐在祂光荣的宝座上时，你们也要坐在十二宝座上，审判以色列十二支派。}\BibleRef{玛 19:28} 对于那些忠信保管了自己塔冷通的人，主也应许了赏报，尽管较小：\BibleQuote{好！你既在少许事上忠信，我必委派你管理许多大事。}\BibleRef{玛 25:21} 因此，信实是当尽的义务，而赏报则在于仁慈。那保持了信实的人，理应得到信实的对待；那获得丰厚利润的人，因不求自己的利益，便赢得了天上赏报的权利。\par

% structure: p0104 source=h2 target=section reason=relative-heading-rank; base=h2

\section{第十三章}

圣安博在讲解福音中关于\OriginalTerm{阉人}{eunuch}的话时，谴责那些使自己成为阉人的人。只有那些通过\OriginalTerm{节欲}{continence}战胜自己的人，才配得赞美；但不应强迫任何人过这种生活，因为基督和宗徒都没有定下这样的法律。因此，婚姻的\OriginalTerm{誓愿}{vow}不应受责备，尽管贞洁的\OriginalTerm{誓愿}{vow}更好。\par

75. 因此，关于此事并没有给出诫命，而是给出了劝谕。贞洁是命令，完全的\OriginalTerm{节欲}{continence}是劝谕。\BibleQuote{这话不是人人所能领悟的，只有那些得了恩赐的人，才能领悟。因为有些阉人，从母胎生来就是阉人，}\BibleRef{玛 25:11} 在这些人身上，存在的是自然的必然性，而非贞洁的德行。\BibleQuote{有些阉人，是自阉的，}是出于自己的意愿，而非出于必然。\BibleQuote{还有些阉人，是被人阉的；……}因此，在他们身上，节欲的恩宠是大的，因为使人节欲的是意志，而非无能。因为完整地保存天主工作的恩赐是合宜的。也不要让他们以为，不受肉身倾向的阻碍是微小的，因为若他们无法获得那奋斗的奖赏，那么罪的课题也被除去了，他们虽不能获得冠冕，却也不再能被战胜。他们还有其他种类的德行，应以此来推荐自己：若是信德坚固，慈悲为怀，远离贪婪，恩宠丰厚。但在他们身上没有罪过，因为他们不知罪的行为。\par

76. 那些\OriginalTerm{自残}{mutilate themselves}的人则不同。我特意提及这一点，因为有些人视这种邪恶的暴力行为为圣事。虽然我不愿对他们发表自己的意见，尽管先辈们的决定仍在；但请想一想，这种行为是否更倾向于显示软弱，而非刚强的声誉。按此原则，便不应战斗以免被战胜，也不应使用双脚以免跌倒，也不应让眼睛发挥功能，因为害怕因\OriginalTerm{情欲}{lust}而跌倒。然而，割损肉身有何益处？因为即使一个眼神也可能有罪？\BibleQuote{凡注视妇女，有意贪恋她的，他已在心里奸淫了她。}\BibleRef{玛 5:28} 同样，凡注视男人，有意贪恋他的，也犯了奸淫。因此，我们应成为贞洁的人，而非懦弱的人；我们的眼睛应是端庄的，而非无力的。\par

77. 因此，没有人应该如许多人所想的那样自阉，而是要获得胜利；因为教会接纳的是得胜者，而非失败者。既然有宗徒的命令在眼前，我又何必使用论证呢？因为你可以找到这样的记载：\Quote{我巴不得那些怂恿你们受割礼的人将自己割绝了。} 因为对于一个生来为得尊荣、装备好去得胜的人来说，为何要失去赢得冠冕和修德的机会呢？他怎能凭着灵魂的勇毅去自阉呢？\BibleQuote{因为有从母胎生来就是阉人的，也有被人阉的，并有为天国的缘故自阉的。} \BibleRef{玛 19:12}\par

78. 然而，这并非给所有人的一条诫命，而是摆在所有人面前的一个心愿。因为发命者必须始终谨守诫命的精确范围，分派任务者必须公平考量；因为：\BibleQuote{两样的法码，为上主所憎恶。} \BibleRef{箴 11:1} 那么，在重量上有过与不及，而教会这两样都不接纳；因为：\BibleQuote{两样的法码，两样的升斗，都为上主所憎恶。} \BibleRef{箴 20:10} 有一些任务，智慧是按各人的德行与力量来估量并分派的。\BibleQuote{能领悟的，就领悟吧！} \BibleRef{玛 19:12}\par

79. 因为万有的创造者知道每个人的性情各不相同，因此用赏报激励德行，而非用锁链束缚软弱。而他，外邦人的导师，我们品行的优良向导，我们内心情感的训诲者，他亲身体验到肉性的法律与心灵的法律相对抗，但屈服于基督的恩宠，我敢说，他知道心灵的各种动态彼此对立；因此，他这样劝勉守贞，既不废除婚姻的恩宠，也不过分高举婚姻而抑制贞洁的愿望。相反，他从推崇贞洁开始，进而提供对治无节制的药方，将崇高召叫的奖赏摆在那较强的人面前，不让任何人在途中昏厥；他赞许领先者，却也不轻视跟随者。因为他自己曾学到，主耶稣赐给一些人大麦饼\BibleRef{若 6:9}，免得他们在途中昏厥；又向另一些人分施祂的圣体\BibleRef{玛 26:26}，好让他们为天国奋力。\par

80. 因为主自己并没有强加这条诫命，而是邀请意志，宗徒也没有立下规条，而是给予了一项劝谕。 \BibleRef{格前 7:25} 但这不是人的劝谕，关于人能力范围内的事，因为他承认天主仁慈的恩赐已赋给了他，好让他知道如何忠实地将以先前者放在首位，并安排后者。因此，他说：\Quote{我认为，}不是，我命令，而是\Quote{我认为这是好的，因为现今的艰难。} \BibleRef{格前 7:26}\par

81. 因此，婚姻的约束不当被视为罪恶而逃避，而是因为它如同一个磨人的重轭而放下。因为法律约束妻子在分娩中受劳苦和疼痛，并要服从丈夫，因为他是她的主。因此，已婚妇女，而非寡妇，在生育儿女时要受劳苦与疼痛；只有已婚者，而非童贞女，在丈夫的权下。童贞女则从这一切中得到释放，她已将自己的感情许给了天主的圣言，她以善意之光照耀的灯盏，等候那应受赞颂的新郎。所以她受劝谕的引导，而不受锁链的束缚。\par

% structure: p0113 source=h2 target=section reason=relative-heading-rank; base=h2

\section{第十四章}

虽然寡妇可能没有领受任何诫命，但她已领受了如此多的劝谕，以致她不该轻视它们。\Person{圣盎博罗削}若给她设下任何圈套，他会感到难过，因为他看见教会的田地因婚姻而更加丰饶；但绝对不能否认，\Person{圣保禄}所赞颂的寡居生活是有益的。因此，他严厉地谈论了那些以法律禁止寡居的人。\par

82. 然而，寡妇同样没有领受任何命令，而是一项劝谕；但这劝谕非只给予一次，而是屡次重复。因为，首先有话说：\BibleQuote{男人不亲近女人，倒好。} \BibleRef{格前 7:1} 再则：\BibleQuote{我愿众人都如同我一样。} \BibleRef{格前 7:7} 又一次：\BibleQuote{他们若常像我一样，为他们倒好。} \BibleRef{格前 7:8} 第四次：\BibleQuote{因现今的艰难，这是好的。} \BibleRef{格前 7:26} 而且，这蒙主悦纳，是尊贵的，末了，恒守寡居更为有福，他不只以此作为自己的判断，也是圣神的愿望。那么，谁能拒绝如此劝导者的善意呢？他将缰绳交给意志，凭自己的经验给他人建议他认为有益之事，他既不轻易被追上，也不因人并驾齐驱而受伤害。那么，谁能退缩，不肯在身体和灵魂上成为圣洁的呢？因为赏报远超辛劳，恩宠胜过所需，工价超过工作。\par

83. 我说这话，并非为给别人设下圈套，而是愿以受托之园地善仆的身份，看见这块教会的田地结实累累，一时绽放纯洁的花朵，一时坚守寡居的庄重，再一时又因婚姻的果实而丰盛。因为它们虽各不同，却同是一块田地的出产；花园中的百合花，不如田野间的麦穗那样多，而且预备接受种子的田地，远比庄稼收割后休耕的土地为多。\par

84. 所以，守寡是好的，宗徒们的判断常常称赞它，因为它是信德的教师，也是贞洁的教师。然而，那些尊崇淫乱和他们神的羞耻之事的人，却对独身和守寡施行惩罚；他们热衷于追求罪行，竟惩罚对德行的追求；名义上为增长人口，实际却是要扼杀贞洁的志向。因为士兵服役期满后，放下武器，离开所守的岗位，以老兵身份退伍返回故土，以便在辛劳一生之后获得休息，并促使他人因指望未来的安息而更甘愿忍受劳苦。同样，农夫年老气衰时，把扶犁之事交给别人，因青春的劳苦而身心俱疲，在老年得以享受事先的积蓄，但仍随时准备修剪葡萄藤，而不是榨取葡萄汁，以抑制早期的疯长，并用修枝刀剪除青春期的放纵，仿佛在教导说，甚至葡萄树也当追求蒙福的丰产。\par

85. 同样地，寡妇犹如老兵，服役期满，虽放下了婚姻生活的武器，却仍安排全家的和平：她虽已免于负重，仍为将要结婚的年轻人警醒守望；她以老年的周到思虑，安排何处更需费力才得效益，何处物产更丰，何者适合作婚姻的联结。所以，如果田地交托给年长者而非年轻者，为何你会认为为人妻比作寡妇更有利呢？但如果信仰的迫害者也成了寡妇的迫害者，那么对持守信仰的人来说，寡妇绝不应被视为惩罚而逃避，而应被视为奖赏而珍视。\par

% structure: p0119 source=h2 target=section reason=relative-heading-rank; base=h2

\section{第十五章}

\Person{圣盎博罗削}回应那些以渴望子女为由而再婚的人，尤其针对那些在前次婚姻中已有子女的人；并指出随之而来的子女间的纷争，甚至夫妻之间的纠纷，并警告勿在这事上错用圣经例证。\par

86. 然而，或许有人以为，为生育子女而再次结婚是好的。但如果渴望子女是结婚的理由，那么已有子女，这理由当然就不存在了。想要再度尝试那已徒然尝试过的生育，或者再度承受那已忍受过的孤寂，难道是明智的吗？这是那些没有子女之人的情况。\par

87. 此外，那些曾生育子女又失去子女的（因为有望生育子女的人会有更强烈的渴盼），我说，难道她不觉得在举行第二次婚礼时，是在遮掩已失子女的死亡吗？她岂不是在重寻痛苦吗？难道她面对希望的坟墓、她所受丧亡的回忆、哀悼者的声音，不会退缩吗？或者，当火炬点燃，黑夜降临时，她岂不更会以为是在预备葬礼，而非洞房？那么，我的女儿，你为何去重寻你所畏惧的忧伤，而不寻求你不再期盼的子女呢？如果忧伤如此沉重，人更应躲避而非追寻那造成忧伤的事。\par

88. 对于已有子女的人，我该给你什么建议呢？你有什么理由再婚呢？也许是愚昧的轻浮，或是不节制的习惯，或是受伤的心灵的意识在催逼你。但劝言是给清醒者的，不是给沉醉者的，所以我的话是对每种情况都健全的自由良心说的。受伤者有疗法，端正者有劝告。那么，我的女儿，你打算做什么呢？你既有自己的子女，为何向外寻求继承人呢？你所渴望的不是子女，因为已有，而是你已获自由的奴役。因为这是真正的奴役，其中爱情枯竭，不再为童贞的魅力、充满圣洁端庄和恩宠的青春年华所激发；冒犯更易感到，粗鲁更被猜疑，和谐更不常见，没有因时间深植的爱或青春鼎盛之美的牢固维系。对丈夫的职责成为重担，以致你害怕爱自己的子女，羞于看他们；而那通常使父母亲情增添的，反而成为纷争的缘由。你愿生育的后代，将不是子女的弟兄，而是他们的仇敌。因为生育其他子女，岂不就是夺去你已有子女的吗？他们因此被剥夺了亲情关怀和财产利益。\par

89. 天主的法律以其权威将夫妻结合在一起，然而彼此相爱仍是难事。因为天主从人身上取了一根肋骨，形成女人，使他们彼此结合，说：\BibleQuote{二人成为一体。}\BibleRef{创 2:24} 他说这话不是指第二婚姻，而是指第一婚姻，因为厄娃没有再嫁，圣教会也不承认第二个新郎。\Quote{这奥秘在基督和教会内是伟大的。\BibleRef{弗 5:32} 同样，依撒格除了黎贝加以外，未曾认识别的妻子，\BibleRef{创 24:67} 也未与撒辣以外的妻子埋葬他的父亲亚巴郎。}\BibleRef{创 25:10}\par

90. 但在圣辣黑耳身上，更多的是奥秘的预像，而非真的婚姻秩序。尽管如此，在她身上也有我们可归于初婚恩宠之处，因为他最爱那最初订婚的人；欺骗并未关闭他的意愿，介入的婚姻也没有摧毁他对未婚妻的爱。因此，圣祖教导我们，该何等重视初婚，因为他自己如此重视最初的婚约。那么，我的女儿，你要留意，免得你既不能持守婚姻的恩宠，又增添自己的烦恼。\par

\SourceRecord{Translated by H. de Romestin, E. de Romestin and H.T.F. Duckworth.；Nicene and Post-Nicene Fathers, Second Series；Vol. 10.；Edited by Philip Schaff and Henry Wace.；Buffalo, NY: Christian Literature Publishing Co.,；1896.；<http://www.newadvent.org/fathers/3408.htm>.}
